\documentclass[10pt,a4paper]{amsart}
\usepackage[margin=19mm]{geometry}
\usepackage{amsmath,amssymb,mathtools}
\usepackage[scr=rsfs,cal=euler]{mathalfa}
\usepackage{xfrac}
\usepackage[unicode,hidelinks,bookmarks=false]{hyperref}
\newcommand{\ie}{i.e.\ }
\newcommand{\cf}{cf.\ }
\newcommand{\resp}{respectively\ }
\newcommand{\Subsection}{\subsection}
\providecommand{\nobreakdash}{\nobreak\hbox{-}}
\setlength{\emergencystretch}{2em}
\multlinegap0pt
\usepackage[shortlabels]{enumitem}
\usepackage{amssymb}
\usepackage{mathtools}
\usepackage{tikz}
\newtheorem{corollary}{Corollary}
\newtheorem{definition}{Definition}
\newtheorem{lemma}{Lemma}
\newtheorem{example}{Example}
\newcommand\Cycle{\mathrm{Cycle}}
\newcommand\NE{\mathrm{NE}}
\newcommand\NW{\mathrm{NW}}
\newcommand\Path{\mathrm{Path}}
\newcommand\SE{\mathrm{SE}}
\newcommand\SW{\mathrm{SW}}
\newcommand\bg{\mathbf{g}}
\newcommand\bk{\mathbf{k}}
\newcommand{\defn}[1]{\emph{\color{blue} #1}} 
\newcommand\east{\mathrm{E}}
\newcommand\north{\mathrm{N}}
\newcommand\south{\mathrm{S}}
\newcommand\squarette{\scalebox{0.75}{\hbox{$\,\square\,$}}}
\newcommand\west{\mathrm{W}}
\title{Locally anti-blocking $\bg$-polytopes for flow polytopes}
\author{Jonah Berggren, Benjamin Braun, Alvaro Cornejo, James Ford McElroy, Chloe' Napier, Zachery Peterson, Williem Rizer, Khrystyna Serhiyenko, Martha Yip}
\date{Source: arXiv:2605.27007v1 (CC BY 4.0)}
\begin{document}
\maketitle

\noindent\emph{Source and licence.} Locally anti-blocking $\bg$-polytopes for flow polytopes. Version arXiv:2605.27007v1 (CC BY 4.0), distributed by arXiv under the Creative Commons Attribution 4.0 International licence, http://creativecommons.org/licenses/by/4.0/, as recorded in the arXiv licence metadata for this exact version and shown on the abstract page https://arxiv.org/abs/2605.27007v1. Full licence notice: \texttt{licenses/arxiv-2605.27007-CC-BY-4.0.txt}.\par\smallskip

\noindent\textit{Editorial scope.} Counted unit: the article's lemma lem:thecoherencelemma (source0.tex lines 1400--1425). The article's complete original proof of it is delivered with it (source0.tex lines 1468--1525, one \texttt{proof} environment of 58 lines, which the article places away from the statement). No mathematics is written, repaired or re-derived: the statement and the proof are the article's own lines, in the article's own notation, and no retained line is substituted or re-flowed.  Retained uncounted as the source-local context the proof needs: lines 1091--1104; lines 1295--1303; lines 1358--1371; lines 1453--1465.  One declared adaptation: exactly 1 ancillary strings inside the retained spans are omitted (image-inclusion commands and, where the source repeats a label, the repeated label definitions) and no other line differs from the source; the figure environments, with their placement, captions and labels, are kept, and the package records the omitted strings in fidelity receipts bound to the affected bodies.  No retained line contains a citation, so the wrapper carries no bibliography and no bibliography entry.  The retained lines contain the article's own diagram code, which the wrapper typesets with the standard tikz packages; no image file is delivered and the package declares no asset dependency.  Editorial formatting only, in addition: an AMS \texttt{amsart} preamble that restates the article's own declarations and macros used by the retained lines, namely the environments \texttt{corollary}, \texttt{definition}, \texttt{lemma}, \texttt{example} on separate counters, together with the standard packages those lines need.  The retained environments print in this document as Corollary 1, Definition 1, Lemma 1, Example 1 in order of appearance, whereas the article prints the same items with the numbering of its own full document.  The complete native source of the article as distributed in the arXiv e-print for this exact version is bundled unchanged as \texttt{sources/201462/source0.tex}.
\begin{corollary}\label{cor:incoherence-square}
    Let $(G,F)$ be a locally anti-blocking framed full DAG, let $P$ be its $\bg$-polytope, and suppose $R$ and $S$ are routes.
    Let $f$ be the minimal face of the $\bg$-polytope containing their $\bg$-vectors. 
    Then:
    \begin{enumerate}
    [label=(\roman*)]
        \item If $E_{\{R,S\}}$ contains the edge set of an exceptional route, then $f$ is the entire polytope $P$ (the routes lie on different facets).
        \item Otherwise:
        \begin{enumerate}
        \item If $R$ and $S$ are in conflict, then $E_{\{R,S\}}$ contains four distinct routes and $f$ is the quadrilateral face that is the convex hull of their $\bg$-vectors.
        \item If $R$ and $S$ are coherent, then $E_{\{R,S\}}$ contains two distinct routes and $f$ is the edge formed by $R$ and $S$.
        \end{enumerate}
    \end{enumerate}
\end{corollary}

\begin{figure}
    \centering
    
    \caption{The coherence diagram corresponding to $\Cycle(3,2)$.
    Boxes with a dot inside indicate exceptional routes.
    Boxes containing a diamond indicate anomalous routes.
    }
    \label{fig:cycle32diagram}
\end{figure}

\begin{definition}\label{def:diagramregions}
    Let $R$ be a route in a locally anti-blocking framed DAG.
    The sets $\north(R)$, $\south(R)$, $\east(R)$, and $\west(R)$ are the regions of the coherence diagram strictly to the north, south, east, and west of $R$ respectively. I.e., $\north(R)$ consists of the boxes above $R$ in its column, $\east(R)$ consists of the boxes to the right of $R$ in its row, and so on.
    In the case of a cycle, when $R$ is in either the top or bottom two rows (which are identified along the lower right and upper left boundary), the regions $\east(R)$ and $\west(R)$ continue across the line of identification.

    We define four additional regions.      \begin{itemize}
            \item The region $\NW(R)$ consists of those boxes $R'$ such that there exists boxes $T_1$ and $T_2$ such that $\{T_1\}=\west(R) \cap \south(R')$ and $\{T_2\}=\north(R) \cap \east(R')$.
            \item The region $\NE(R)$ consists of those boxes $R'$ such that there exist boxes $T_1$ and $T_2$ such that $\{T_1\}=\east(R) \cap \south(R')$ and $\{T_2\}=\north(R) \cap \west(R')$.
            \item The region $\SW(R)$ consists of those boxes $R'$ such that there exist boxes $T_1$ and $T_2$ such that $\{T_1\}=\west(R) \cap \north(R')$ and $\{T_2\}=\south(R) \cap \east(R')$.
            \item The region $\SE(R)$ consists of those boxes $R'$ such that there exist boxes $T_1$ and $T_2$ such that $\{T_1\}=\east(R) \cap \north(R')$ and $\{T_2\}=\south(R) \cap \west(R')$.
        \end{itemize}
    We say that a route $S$ is \defn{northwest} of $R$ if $S\in \NW(R)$, and similarly for northeast, southwest, and southeast.
    Supposing $R$ is northwest of $S$, the \defn{2-by-2 array} subtended by $R$ and $S$, denoted $R\squarette S$, is the set $\{R, S, T_1, T_2\}$ where $T_1$ and $T_2$ are the routes required to define $R$ as northwest of $S$. Thus the notation $R\squarette  S$ implies $R\in \NW(S)$.
\end{definition}

\begin{example}\label{ex:anomalouslemma}
    Consider $\Cycle(3,2)$ with coherence diagram given in Figure~\ref{fig:cycle32diagram}.
    The anomalous routes are indicated by diamonds.
    The route $(1^-,4^-)$ is northwest of $(1^+,4^+)$, $(4^+,1^-)$, and $(4^+,1^+)$, and it is straightforward to check that these three routes are the only anomalous routes in conflict with $(1^-,4^-)$.

    The route $(1^+,4^-)$ is southeast of $(4^-,1^-)$, due to the identified edges on the diagram, and it is northwest of $(4^+,1^-)$ and $(4^+,1^+)$.
    It is straightforward to check that these are the only anomalous routes in conflict with $(1^+,4^-)$.

    The route $(1^+,4^+)$ is southeast of $(1^-,4^-)$, $(4^-,1^-)$, and $(4^+,1^-)$, due to the identified edges on the diagram.
    One can check that these are the only conflicting anomalous routes with $(1^+,4^+)$.

     Further, it is straightforward to check that the union of edges in any of these conflicting pairs contains an exceptional route, and hence by Corollary~\ref{cor:incoherence-square} the minimal face containing that pair is $P$.
\end{example}

\begin{lemma}[The Coherence Lemma]\label{lem:thecoherencelemma}
    Let $R$ and $S$ be routes in a locally anti-blocking framed DAG $G$ with $\bg$-polytope $P$.
    \begin{enumerate}[label=(\roman*)]

    \item If both $R$ and $S$ are anomalous:
    \begin{enumerate}
        \item $S\in \NW(R)\cup \SE(R)$ if and only if $R$ and $S$ are in conflict, and in this case the entire polytope $P$ is the minimal face containing both $R$ and $S$. 
        \item If $R$ and $S$ are coherent:
        \begin{enumerate}
            \item If one if $R$ or $S$ is exceptional, the minimal face containing them is $P$.
            \item Otherwise, the minimal face containing them is the edge from $R$ to $S$.
        \end{enumerate}
    \end{enumerate}
    \item If $R$ and $S$ are not both anomalous:
    \begin{enumerate}
        \item $S\in \NW(R) \cup \SE(R)$ if and only if $S$ is in conflict with $R$.
        Further, setting $R\squarette  S =: \{R, S, T_1, T_2\}$ as in Definition~\ref{def:diagramregions}, we have:
    \begin{enumerate}
        \item if one of $\{T_1,T_2\}$ is exceptional, then $R$ and $S$ lie on different facets of $P$.
        \item Otherwise, the convex hull of the routes in $R\squarette  S$ forms the minimal face of $P$ containing $R$ and $S$.
        \end{enumerate}     
        \item $S \in \NE(R) \cup \SW(R)$ if and only if $S$ is coherent with $R$ and the line segment from $R$ to $S$ is not an edge of $P$.
        \item $S$ lies outside of $\NW(R)\cup \NE(R)\cup \SE(R)\cup \SW(R)$ if and only if $S$ is coherent with $R$ and the line segment from $R$ to $S$ is an edge of $P$.
    \end{enumerate}
    \end{enumerate}
\end{lemma}

\begin{proof}[Proof of Lemma~\ref{lem:thecoherencelemma}]

The proof of \textit{(i)} is a tedious case-by-case verification.
Note that there are only eight anomalous routes in any $\Cycle(k_1,k_2)$, and two of them are exceptional.
These cases are an extension of the partial analysis of pairs of anomalous routes given in Example~\ref{ex:anomalouslemma}.

We now consider the proof of \textit{(ii)}, which is the general case.
We begin by proving \textit{(ii)(a)}, which states $S\in \NW(R) \cup \SE(R)$ if and only if $S$ is in conflict with $R$.
Observe that once we prove \textit{(ii)(a)}, the claims in \textit{(ii)(a)(i)} and \textit{(ii)(a)(ii)} follow immediately from Corollary~\ref{cor:incoherence-square}.
For the forward implication, without loss of generality suppose that $S\in \NW(R)$.
The routes in $R\squarette  S$ correspond to the bounding corners of a rectangle within the coherence diagram.
We argue by cases as follows.
\begin{enumerate}
    \item Suppose three things: $R\squarette S$ is contained in the bounding strip, vertex $i$ is adjacent to the source, and $S$ contains $i^-$ and $R$ contains $i^+$ (the case where $i$ is adjacent to the sink is a similar argument).
    One of $i^+$ or $i^-$ might be denoted by $i$ in the case that $G=\Path(\bk)$ and $i$ is a leaf of the inner graph. 
    If the edges connecting $S$ and $R$ to the source are on different legs, then the routes conflict at $i$.
    Otherwise, since $R$ is to the right of $S$ in the coherent diagram, either $R$ exits from the leg containing $R\cap S$ on a different vertex from the one where $S$ exits or they both exit the leg at the final possible vertex.
    If they exit on different vertices and if $R\cap S$ lies on a $2$-leg, then $R$ exits the leg after $S$.
    If they exit on different vertices and if $R\cap S$ lies on a $1$-leg, then $S$ exits the leg after $R$.
    In both cases, $R$ and $S$ are in conflict.
    If $R$ and $S$ both enter and exit the leg at the source- and sink-adjacent vertices, respectively, then (recalling that $R$ and $S$ are not both anomalous) $R\squarette S$ must form a contiguous $2\times 2$ square in a corner of the bounding strip.
    Hence, $R\cap S$ is an entire leg and we have that $R$ and $S$ conflict across the leg.
    \item Suppose that $R\squarette S$ is not contained in the bounding strip, i.e., at least one element of $R\squarette S$ is contained in a staircase shape outside of the bounding strip, so that element corresponds to a route passing through the inner vertices of a fixed leg $L$ in $G$.
    Every element of $R\squarette S$ therefore is contained in the portion of the coherence diagram given by the routes that lie on $L$.
    Thus, the first and last inner vertices in the route $(R\cup S)\cap L$ are distinct, and since $S\in \NW(R)$, one of these is contained in $R$ and the other is contained in $S$.
    Since $R$ and $S$ lie on the same leg, and since the framing labels are constant for the edges on the leg, it must be that $R$ and $S$ conflict across $R\cap S$.
\end{enumerate}

We next prove the reverse implication for \textit{(ii)(a)}.
Suppose that $R$ and $S$ are in conflict, where $R$ and $S$ are not both anomalous.
The connected path $R\cap S$ must be contained within a leg of $G$, since the inner graph of $G$ is either a path or a cycle.
We consider the following cases.
\begin{enumerate}
    \item Suppose that for the conflicting routes $R$ and $S$, there exists a vertex $i$ with two edges connected to either the source or the sink such that one of the routes contains $i^+$ while the other contains $i^-$.
    Note that in the case where $G=\Path(\bk)$, one of these edges might be denoted by $i$ if the vertex $i$ is a leaf of the inner graph.
    We will assume that $i$ is adjacent to the source, as the argument is similar when $i$ is adjacent to the sink.
    Thus, $R$ and $S$ are both contained in the bounding strip for the diagram, within the rectangle indexed by rows $i^+$ and $i^-$.
    Further, $R$ and $S$ are in different rows and columns, since they must have distinct sink-adjacent edges, otherwise the routes would be identical to the right of vertex $i$.
    The other two routes contained in $R\cup S$ are precisely the routes given by the corners of the rectangle in the diagram defined by $R$ and $S$, as the route pairs for those two routes are obtained by exchanging the sink-adjacent edges in the route pairs for $R$ and $S$.
    Hence, either $R\in \NW(S)$ or $S\in \NW(R)$.
    \item Suppose that $R\cup S$ does not contain both $i^+$ and $i^-$ for any source-adjacent or sink-adjacent vertex $i$.
    In this case, the source-adjacent vertices for $R$ and $S$ are distinct, as are the sink-adjacent edges.
    Thus, the route pairs for $R$ and $S$ have distinct first coordinates and distinct second coordinates, and the other two routes in $R\cup S$ are obtained by exchanging the second coordinates in the route pairs for $R$ and $S$.
    This exchange of second coordinates produces routes corresponding to the $T_1$ and $T_2$ in $R\squarette S$, and hence either $R\in \NW(S)$ or $S\in \NW(R)$.
\end{enumerate}

We next prove cases \textit{(ii)(b)} and \textit{(ii)(c)}. 
Observe that due to case \textit{(ii)(a)}, it is immediate that $R$ and $S$ are coherent if and only if $S\notin \NW(R)\cup \SE(R)$.
This establishes the ``coherence'' condition in the if and only if statements for \textit{(ii)(b)} and \textit{(ii)(c)}.
We need to verify that $S\in \NE(R)\cup\SW(R)$ precisely when the line segment from $R$ to $S$ is not an edge of $P$.
Without loss of generality, suppose $S\in \NE(R)$.
By definition of $\NE(R)$, we have $S\in \NE(R)$ if and only if there exist conflicting routes $T$ and $U$ such that $T\squarette U=\{T,U,R,S\}$.
By part \textit{(ii)(a)} of Corollary~\ref{cor:incoherence-square}, it follows that the minimal face containing $T\squarette U$ is the quadrilateral face that is the convex hull of their $\bg$-vectors.
Because the quadrilateral is the smallest face containing $T$ and $U$, the line segment from $R$ to $S$ passes through the interior of this quadrilateral, hence is not an edge of $P$.

As a consequence of \textit{(ii)(a)} and \textit{(ii)(b)}, we have that routes outside of $\NW(R)\cup \NE(R)\cup \SE(R)\cup \SW(R)$ are precisely those that are coherent and for which the line segment in $P$ between them forms an edge of $P$.
This completes the proof of \textit{(ii)(c)}, and thus the proof of the lemma is complete.
\end{proof}

\end{document}
